A eq. (3.13) expressa a velocidade $v$ em termos de propriedades musculares e corporais de animais geometricamente semelhantes. Confirma-se que $v$ é independente do comprimento característico $L$ pelo seguinte argumento de escala energética.

\textbf{Hipóteses de semelhança geométrica:}
\begin{itemize}
\item massa corporal: $m \propto \rho_0\,L^3$ (densidade média $\rho_0$ constante);
\item força muscular máxima: $F_{\max} \propto \sigma_0\,L^2$ (tensão muscular $\sigma_0$ constante);
\item comprimento do passo: $s \propto L$.
\end{itemize}

\textbf{Balanço energético por passada:}
O trabalho mecânico disponível por passo é da ordem de força vezes distância:
\[
W_{\rm paso} \sim F_{\max}\cdot s \sim \sigma_0\,L^2 \cdot L = \sigma_0\,L^3.
\]
Esse trabalho é convertido em energia cinética máxima:
\[
E_k = \tfrac{1}{2}\,m\,v^2 \sim \rho_0\,L^3\,v^2.
\]

\textbf{Equação de escala:}
Igualando $W_{\rm paso} \sim E_k$:
\[
\sigma_0\,L^3 \sim \rho_0\,L^3\,v^2.
\]
O fator $L^3$ cancela em ambos os lados, resultando em:
\[
v^2 \sim \frac{\sigma_0}{\rho_0} = \text{constante},\qquad\text{ou seja,}\qquad v \sim \sqrt{\frac{\sigma_0}{\rho_0}} \propto L^0.
\]

\textbf{Conclusão:} A eq. (3.13) mostra que $v \propto L^0$: a velocidade máxima de corrida \textbf{não depende do tamanho do animal}, dependendo apenas das propriedades mecânicas internas dos músculos ($\sigma_0$ e $\rho_0$).

